\documentclass[11pt,a4paper]{article}
\usepackage[margin=0.95in]{geometry}
\usepackage{booktabs}
\usepackage{amsmath}
\usepackage{graphicx}
\usepackage{xcolor}
\usepackage{titlesec}
\usepackage{enumitem}
\usepackage[hidelinks]{hyperref}
\usepackage{fancyhdr}
\usepackage{caption}

\definecolor{ink}{HTML}{1A1A1A}
\definecolor{muted}{HTML}{5A5A5A}
\definecolor{rule}{HTML}{BBBBBB}
\color{ink}

\titleformat{\section}{\normalfont\large\bfseries}{\thesection.}{0.5em}{}
\titleformat{\subsection}{\normalfont\normalsize\bfseries}{\thesubsection}{0.5em}{}
\titlespacing*{\section}{0pt}{14pt}{5pt}
\titlespacing*{\subsection}{0pt}{10pt}{3pt}

\captionsetup{font=small,labelfont=bf,skip=4pt}

\pagestyle{fancy}
\fancyhf{}
\fancyhead[L]{\footnotesize\color{muted}NeuFlow v3}
\fancyhead[R]{\footnotesize\color{muted}\thepage}
\renewcommand{\headrulewidth}{0.2pt}

\setlist[itemize]{topsep=3pt,itemsep=2pt,parsep=0pt,leftmargin=1.3em}
\setlist[enumerate]{topsep=3pt,itemsep=2pt,parsep=0pt,leftmargin=1.5em}
\setlength{\parskip}{5pt}
\setlength{\parindent}{0pt}

\newcommand{\px}{\,px}
\newcommand{\ms}{\,ms}

\begin{document}

\begin{center}
{\LARGE\bfseries NeuFlow v3}\\[4pt]
{\large Queried optical flow with calibrated uncertainty\\for region-directed perception}\\[10pt]
{\normalsize Shriman Raghav Srinivasan}\\[2pt]
{\footnotesize\color{muted} MS Robotics, Northeastern University \quad$\cdot$\quad Field Robotics Lab \quad$\cdot$\quad August 2026}
\end{center}
\vspace{2pt}
{\color{rule}\hrule}
\vspace{8pt}

\section*{Abstract}
Real-time optical flow networks emit a displacement for every pixel, at one fixed
resolution, on every call, and retain no state between calls. Many consumers of
flow, including registration, sparse tracking, visual odometry and survey
mosaicking, need a few hundred correspondences at locations they select, and a
platform working to a millisecond budget can afford neither the full field nor
the recomputation. This report presents NeuFlow v3, which retains the NeuFlow v2
pipeline through its coarse flow stage and replaces only the final learned convex
upsampler with an implicit decoder evaluated at arbitrary continuous coordinates.
The front end is frozen and verified bit-identical across all 137 shared tensors
after every run, so every measured difference is attributable to the decoder
alone. The decoder samples four feature sources in a $3\times3$ window around a
query, fuses them through a gated network, and emits a convex combination of ten
candidate displacements drawn from the coarse field, bounding the prediction
inside the locally supported motion range, together with a per-query error scale
trained under a Laplace likelihood.

The resulting factorisation into a cached coarse pass and a queried decode changes
the cost of $N$ queries about one frame pair from $33.3N$\ms\ to $32.8 + 1.25N$\ms.
A repeat query on a cached pair costs 1.25\ms\ against 33.3\ms, and an object
appearing in a frame already being processed is answered from cache in 1.68\ms\ at
2.10\px, against 7.29\ms\ and 3.60\px\ for a fresh crop through the baseline, which
is both slower and less accurate because cropping discards the context global
matching requires. Sparse output reproduces the dense field to 0.00057\px, so
sparse and dense decoding are one function evaluated on two point sets rather than
two operating modes. Across 2{,}348{,}000 held-out queries the predicted error
scale orders measured error monotonically over five equal-count bins, from
0.48\px\ to 7.10\px; abstaining on the least confident fifth reduces mean error
over the accepted 80\% to 1.06\px, a selection the baseline cannot make at any
threshold because it emits flow alone.

The cost is quantified rather than avoided. Mean end-point error rises from
2.324\px\ to 2.384\px, 2.6\%, from a model 13\% smaller, and the like-for-like
single-dataset comparison is 2.500\px. The penalty is concentrated in sub-pixel
precision, 1.5 points of 1-pixel accuracy, rather than in large errors, and dense
reconstruction is 11\% slower, so the baseline remains the correct choice for
consumers that take one dense field per frame. Because every driving result is
measured on a domain v3 trained on and the baseline did not, an unconfounded
cross-domain comparison is reported on aerial imagery neither model trained for,
where the queried decoder records lower error at all seven crop settings tested,
by consistent but small margins. Finally, region-directed inference for
fast-moving platforms is characterised: processing 14\% of the frame costs
0.034\px, and the context margin a cropped region requires is set by inter-frame
displacement rather than by image size, a relation that collapses across a
threefold difference in motion scale and can therefore be sized from platform
speed and frame rate before any flow is computed.

\section{Why queryable flow}
Flow networks produce a dense map. Many consumers do not want one.

\begin{itemize}
\item Registration needs a few hundred correspondences at points it picks itself.
\item Sparse tracking needs flow only at feature points.
\item Survey mosaicking needs matches in overlap regions, often between pixels.
\end{itemize}

Computing 479{,}232 values to use 800 of them is the difference between real time
and not on constrained hardware. A fixed-resolution design also cannot answer any
query above its input sampling rate, which is what high-resolution registration
needs.

The question this project asks: can the last stage be made queryable without
losing accuracy or speed?

\section{Related work}
\label{sec:related}

Flow started as a variational problem. Horn and Schunck~\cite{hs81} minimise a
global energy pairing brightness constancy with a smoothness penalty on the
field. Lucas and Kanade~\cite{lk81} solve brightness constancy in closed form
over a window, assuming the flow is constant inside it. Both are local by
construction, because the data term is a first-order expansion about zero
displacement and holds only while motion is small relative to image structure.
Brox et al.~\cite{brox04} add gradient constancy and a warping scheme that solves
at reduced resolution first, which extends the usable range but still loses a
small object that moves far enough to vanish from the coarsest level. Large
displacement of small structures is the failure every learned method since has
been measured against.

FlowNet~\cite{flownet} replaced the energy with a network trained end to end on
synthetic data, and showed that an explicit correlation layer between the two
frames helps more than letting a generic encoder discover matching on its own. It
was well short of variational accuracy. FlowNet 2.0~\cite{flownet2} closed that
gap by stacking several such networks with warping between them and a separate
branch for small displacement, at the cost of a large model.

PWC-Net~\cite{pwcnet} and LiteFlowNet~\cite{liteflownet} put the classical
pyramid back inside the network. Features are extracted at several scales, flow
from the coarser level warps the second frame, and a cost volume is built over a
bounded displacement range around that warp. Bounding the search keeps the volume
small, which is what made these models compact enough to be practical. The cost
is the cascade's own weakness: an error made at a coarse level is never revisited.
RAFT~\cite{raft} removed the cascade by building one all-pairs correlation volume
at $1/8$ resolution and reading from it repeatedly with a recurrent unit,
refining a single estimate at fixed resolution. It set the accuracy standard the
field still works against, and it is slow. GMFlow~\cite{gmflow} treats
correspondence as global matching instead, with cross-attention over the whole
feature map producing flow directly and refinement reduced to a correction.

Real-time work asks what of this survives a millisecond budget.
NeuFlow~\cite{neuflow1} keeps global matching but runs it at $1/16$ behind a
shallow backbone, then refines with cheap recurrent updates.
NeuFlow v2~\cite{neuflow2} restructures the backbone and refinement and ends with
a learned convex upsampler that predicts blending weights over a $3\times3$
neighbourhood of coarse flow. That upsampler is the stage v3 replaces, and
everything above it is what v3 inherits frozen.

Implicit representations arrived from elsewhere. NeRF~\cite{nerf} showed a
coordinate-conditioned MLP can hold a scene and be evaluated anywhere in it, and
SIREN~\cite{siren} showed the activation function decides how much
high-frequency detail such a network can hold at all. LIIF~\cite{liif} made the
idea discrete to continuous for images: a local latent code plus a relative
coordinate, decoded by one shared MLP, so a single trained model can be sampled
at any resolution. AnyFlow~\cite{anyflow} brought arbitrary-scale decoding to
flow and predicts convex weights over candidate displacements, which is where
v3's ten-candidate head comes from. InfiniDepth~\cite{infinidepth} does the
continuous version for depth and contributes the gated hierarchical fusion of
multi-scale features that v3's decoder uses.

Confidence has its own line. Kendall and Gal~\cite{kendall} separate aleatoric
from epistemic uncertainty and train a heteroscedastic likelihood whose predicted
scale is supervised only by the error it explains. Ilg et al.~\cite{ilg18}
applied this to flow directly, comparing multi-hypothesis networks against a
single predictive distribution and showing the estimate is usable rather than
decorative. PDC-Net~\cite{pdcnet} predicts a mixture density per correspondence
and uses it to decide which matches to trust, which is the same use v3 makes of
its per-query scale.

Every method above emits a dense field at one fixed resolution, chosen when the
network is built. The queryable ones are queryable in scale, not in coordinate,
and the implicit decoders that are queryable in coordinate have not been asked to
run inside a network with a millisecond budget. So whether arbitrary-coordinate
querying is affordable at speed is open. That is what this report measures.

\section{Method}

\subsection{What is inherited}
NeuFlow v2~\cite{neuflow2} runs a shallow CNN
backbone at $1/8$ and $1/16$ resolution, cross-attention and global matching at
$1/16$ for an initial flow, recurrent refinement (one iteration at $1/16$, eight
at $1/8$), then a learned convex upsampler to full resolution.

v3 keeps all of that up to the $1/8$ coarse flow, with every learned weight
frozen and every batch-normalisation running statistic held at v2's value. I
verify this per checkpoint: all 137 shared tensors are bit-identical to v2, so the
comparison below isolates the decoder and nothing else.

\subsection{The decoder}
Two phases.

\textbf{Phase 1, once per frame pair, 32.8\ms.} The frozen pipeline produces the
coarse flow and feature maps. These are cached.

\textbf{Phase 2, once per query batch, 1.25\ms\ for up to 2{,}048 points.} For a
coordinate $(x, y)$ the decoder samples $3\times3$ windows from four sources
($1/8$ context, $1/8$ features, $1/16$ features, flow-warped frame-1 features),
fuses them in a gated MLP, and predicts weights over ten candidates: the nine
coarse-flow values around the query plus a bilinear sample.

Output is a convex combination of those candidates. That bounds the prediction
inside the locally supported motion range, so the decoder cannot invent
displacement its inputs do not support. The head is initialised so the softmax
sits on the bilinear candidate, which makes the untrained decoder exactly
equivalent to bilinear upsampling (measured: 0.011\px\ maximum deviation).
Training can only move away from that by learning.

\subsection{Uncertainty head}
An optional eleventh output predicts a per-query error scale $b$ under a Laplace
likelihood. It costs no measurable inference time and attaches a confidence value
to every correspondence returned.

\subsection{Interface}
\begin{small}
\begin{verbatim}
state = model.infer_coarse_state(img0, img1)          # once per pair
flow  = model.decode_queries(state, query_coords=q)   # [B,N,2] -> [B,N,2]
flow  = model.decode_queries(state, target_h=H, target_w=W)
flow, b = model.decode_queries(state, query_coords=q, return_uncertainty=True)
\end{verbatim}
\end{small}

Coordinates are continuous. $(312.7, 188.2)$ is as valid as $(312, 188)$. Sparse
queries reproduce the dense field exactly at matching coordinates, to
0.00057\px.

\section{Setup}
\textbf{Evaluation.} VKITTI2 Scene18 and Scene20: 1{,}174 pairs, 460{,}573{,}660
valid pixels, per-pixel metrics. Both scenes are excluded from every training set.
A guard in the loader enforces it and I verify no overlap at pair and frame level
before each run.

\textbf{Controls.} Every training run is generated from one template, so any two
differ in exactly one variable. All use seed 1234, batch 16, learning rate
$2\times10^{-4}$ on a OneCycle schedule, $\gamma = 0.8$, 4{,}096 supervision
queries per image, refinement of $1\times$s16 and $8\times$s8 matching evaluation,
and 100{,}000 steps.

\textbf{Pre-flight.} A ten-check suite runs on CPU and GPU before any job:
batch-normalisation frozen, only decoder parameters trainable, zero-initialised
decoder equals bilinear, sparse equals dense, uncertainty present in both decode
paths, stride-2 fidelity, seed reproducibility.

\textbf{Hardware.} RTX 4060 laptop, fp16, $384\times1248$ input unless stated.

\section{Results}
\label{sec:results}

\subsection{Accuracy}\label{sec:accuracy}
\begin{center}
\small
\begin{tabular}{llrrrr}
\toprule
Model & Training data & EPE (px) & 1px (\%) & 3px (\%) & Params \\
\midrule
NeuFlow v2 & FlyingThings & \textbf{2.324} & \textbf{77.63} & \textbf{89.80} & 9.03\,M \\
NeuFlow v3 & FlyingChairs & 2.500 & 72.81 & 87.88 & \textbf{7.83\,M} \\
NeuFlow v3 & \quad + VKITTI2 & 2.398 & 75.74 & 88.94 & 7.83\,M \\
NeuFlow v3 & \quad + MPI-Sintel & 2.392 & 75.83 & 88.98 & 7.83\,M \\
NeuFlow v3 & \quad + uncertainty head & 2.384 & 76.13 & 89.02 & 7.83\,M \\
\bottomrule
\end{tabular}
\end{center}

Every v3 configuration sits above v2. The best, at 2.384\px, is 2.6\% worse on
mean error and 1.5 points worse on 1-pixel accuracy. The like-for-like comparison
is worse still: trained on FlyingChairs alone, mirroring v2's own training, v3
reaches 2.500\px, which is 7.6\% behind.

So the implicit decoder is less accurate than v2's convex upsampler. The reason is
identified in Section~\ref{sec:limitations} and is not a training artefact: the
decoder's finest input sits at $1/8$ resolution while v2's upsampler reads the
full-resolution frame.

The ordering is monotonic. Each addition of data helps, and the uncertainty head is
best of the four, though the last three differ by less than 0.02\px\ and I would
not read much into their ranking from a single seed.

One qualification applies to this table and to every driving result that follows.
v3's training mix includes other VKITTI2 scenes while v2's does not, so the
comparison is biased in v3's favour on this domain. The accuracy gap above is
therefore a lower bound on the true cost of the decoder, and the only
unconfounded comparison in this report is the aerial one in Section~\ref{sec:aerial}, on a
domain neither model trained for.

\subsection{Compute}
\begin{center}
\small
\begin{tabular}{lrl}
\toprule
Mode & Latency & Note \\
\midrule
NeuFlow v2, full frame & 33.3\ms & its only mode \\
v3 sparse, first query on a new pair & 34.1\ms & level with v2 \\
\textbf{v3 sparse, repeat query on a cached pair} & \textbf{1.25\ms} & \textbf{27$\times$ cheaper} \\
v3 dense, stride 2 & 37.1\ms & not the intended mode \\
\bottomrule
\end{tabular}
\end{center}

The sparse path is a 32.8\ms\ coarse pass inherited from v2 plus a 1.25\ms\
decode. Global matching needs whole-image context, so the coarse pass cannot be
skipped. It is 96\% of a first query, which is why a first query costs what a v2
frame costs.

State reuse is the property that matters. v2 keeps nothing between calls, so
asking a second question about the same frame means recomputing everything. v3
answers from cached state in 1.25\ms. For anything that queries a frame more than
once (iterative registration, RANSAC refinement, interactive selection,
multi-hypothesis tracking) that is a structural difference, not a margin.

Decode latency is 1.254\ms\ at $N=800$ and 1.285\ms\ at $N=2{,}048$. The decode is
bound by kernel launch overhead, not compute, so 2{,}048 points cost the same as
800.

\subsection{Calibrated uncertainty}
\begin{center}
\small
\begin{tabular}{lr}
\toprule
Predicted scale $b$ & Actual mean error \\
\midrule
0.01 to 0.10 & 0.480\px \\
0.10 to 0.19 & 0.896\px \\
0.19 to 0.39 & 1.018\px \\
0.39 to 1.22 & 1.837\px \\
above 1.22 & 7.100\px \\
\bottomrule
\end{tabular}
\end{center}

Measured over 2{,}348{,}000 queries. Actual error rises monotonically across all
five bins, a 15$\times$ span, Pearson $r = 0.318$. The correlation is not strong
but the ordering is reliable and directly usable: weight correspondences
in RANSAC, drop unreliable ones, or send more queries where the model is unsure.
The ordering can be trusted; the per-query magnitude cannot, and I do not claim it.
v2 outputs flow only and has nothing comparable.

\subsection{Abstention: what the ordering is worth}
Because the model can rank its own answers it can decline to answer some of them.
Sorting the same 2{,}348{,}000 queries by predicted confidence and answering only
the most confident fraction gives an accuracy-coverage curve.

\begin{center}
\small
\begin{tabular}{lrr}
\toprule
Coverage & EPE of the accepted set & v2 \\
\midrule
20\% & \textbf{0.480\px} & no selection possible \\
80\% & \textbf{1.058\px} & no selection possible \\
100\% & 2.266\px & 2.324\px \\
\bottomrule
\end{tabular}
\end{center}

Discarding the least confident fifth reduces error over the remaining 80\% of the
frame to 1.058\px. The claim I make without qualification is the internal one:
across 2.35\,M of my own queries, error falls monotonically as predicted
confidence rises, and the spread from the most to the least confident group is
more than fourfold. That is measured on my own outputs and depends on no
comparison.

The comparison against v2 is the weaker framing and I would rather state its
weakness than have it extracted. The accepted 80\% is the easier 80\%, and v2
would also score better on that subset than on the full frame. What survives
regardless is the capability: there is no threshold that can be set on a model
emitting flow alone to produce this curve at all. v2 has one operating point,
structurally. If a query is rejected the consumer asks elsewhere or falls back on
its previous estimate, and asking again costs about a millisecond.

\subsection{Aerial domain: the one unconfounded comparison}\label{sec:aerial}
Every driving result above is measured on a domain v3 trained on and v2 did not.
TartanAir is not: neither model has trained on aerial imagery, which makes it the
only clean head-to-head in this study. The same 40 held-out pairs are scored at
the full frame and at six crop margins.

\begin{center}
\small
\begin{tabular}{lrr}
\toprule
Crop margin & NeuFlow v2 & NeuFlow v3 \\
\midrule
Full frame & 0.781 & \textbf{0.777} \\
0\,px & 1.205 & \textbf{1.176} \\
8\,px & 0.901 & \textbf{0.881} \\
16\,px & 0.872 & \textbf{0.865} \\
24\,px & 0.879 & \textbf{0.854} \\
32\,px & 0.834 & \textbf{0.814} \\
64\,px & 0.793 & \textbf{0.778} \\
\bottomrule
\end{tabular}
\end{center}

v3 records the lower mean end-point error at every one of the seven settings. The
differences are small and worth sizing out loud: they run from 0.004 to 0.029\px,
on 40 pairs and a single seed, so no individual difference is meaningful. What
makes the table worth showing is that the direction repeats across seven
independent settings.

My explanation is a hypothesis rather than a finding: v3 trained on three datasets
where v2 trained on one, and a bounded decoder degrades more gracefully on
unfamiliar input because it cannot leave the locally supported motion range
whatever the network predicts. It is also not uniformly better. On large-motion
failures v3 is worse, 7.2\% against 6.5\%. Forty pairs on one seed is a consistent
signal, not a reversal of the accuracy result in Section~\ref{sec:accuracy}.

\subsection{Interactive tool}
A PyQt application loads a video, steps frame by frame, and computes flow for a
region the user drags out. Two modes: exact decoding inside the box against a
full-frame coarse pass, or a cropped mode running the whole pipeline on the
selection so cost scales with the area requested. Both print their latency
breakdown live.

\section{Flowing only part of the frame}
A fast platform cannot afford full-frame flow every frame, so it flows only the
region that matters. Three questions follow: what cropping to a region costs,
when it pays, and what happens when a new object appears in a frame already being
processed.

\subsection{What a region costs}
\begin{center}
\small
\begin{tabular}{lrrrr}
\toprule
Region processed & Area & Latency & Speedup & EPE in region \\
\midrule
Full frame & 100\% & 33.3\ms & 1.0$\times$ & 0.657\px \\
Region, no margin & 7.9\% & 7.8\ms & 4.3$\times$ & 1.089\px \\
\textbf{Region + 32\,px margin} & \textbf{13.8\%} & \textbf{7.6\ms} & \textbf{4.4$\times$} & \textbf{0.691\px} \\
Region + 64\,px margin & 20.1\% & 8.1\ms & 4.1$\times$ & 0.667\px \\
Region + 128\,px margin & 34.2\% & 9.9\ms & 3.4$\times$ & 0.655\px \\
\bottomrule
\end{tabular}
\end{center}

Keeping full resolution and processing 14\% of the frame costs 0.034\px\ and runs
4.4$\times$ faster. Beyond a 32\px\ margin nothing improves, so the extra area is
wasted.

Cropping applies to any flow network, v2 included. It is a platform technique,
not a property of the queryable decoder.

\subsection{The margin rule}
A crop removes the context global matching uses to find large displacement. With
no margin, error rises 0.432\px\ and 24\% of large-motion pixels fail outright. A
32\px\ margin on 26.6\px\ mean motion costs 0.034\px\ and is the knee; beyond it
nothing improves. The rule is \textbf{margin $\approx$ expected inter-frame
motion $\approx$ speed / frame rate}, so a faster platform needs a larger margin
and saves proportionally less.

\subsection{The rule across two motion scales}
Tested on aerial sequences, which move very differently from driving: 9.26\px\
mean in-region displacement against 26.6\px.

\begin{center}
\small
\begin{tabular}{lrrr}
\toprule
Domain & Mean motion & Margin needed & Penalty there \\
\midrule
Driving & 26.6\px & $\sim$32\px & +0.034\px \\
Aerial & 9.26\px & $\sim$8\px & +0.120\px \\
\bottomrule
\end{tabular}
\end{center}

Both begin at the same 0.43\px\ penalty with no margin, and both have shed most
of it once the margin reaches one frame of motion. The requirement is therefore
set by displacement rather than by image size, and can be sized from platform
speed and frame rate before any flow is computed. The agreement is close up to
that point; the residual beyond it differs by scene, so this predicts the scale
of the margin rather than the exact curve.

\subsection{A new object mid-frame}
\begin{center}
\small
\begin{tabular}{lrr}
\toprule
Policy & Cost of the new object & EPE on it \\
\midrule
v3, sparse queries (800 pts) & \textbf{1.68\ms} & 2.10\px \\
v3, dense ROI & 4.36\ms & 1.77\px \\
v2, new crop & 7.29\ms & 3.60\px \\
v3, new crop & 8.23\ms & 3.61\px \\
\bottomrule
\end{tabular}
\end{center}

Because the coarse pass is cached, the decoder answers a newly appeared object
for 1.68\ms. A cropped pipeline must run an entire new pass and is markedly less
accurate doing so, having lost the global context. NeuFlow v2 keeps no state
between calls and has nothing to answer from. This is the one result in this
section specific to the architecture rather than to cropping: answering a
question that was not anticipated when the frame arrived.

Two disjoint crops cost more than one full frame, so the policy rule is to crop
once or not at all.

\section{What v3 is for}
The contribution is an operating point, not a leaderboard position. v3 trades
2.6\% of mean accuracy and 1.5 points of 1-pixel accuracy for three properties the
fixed-resolution design cannot express, from a model 13\% smaller:

\begin{enumerate}
\item Flow at any continuous coordinate, with sparse output identical to dense.
\item Repeat access at 1.25\ms\ per query batch against a cached frame, versus a
      full 33.3\ms\ recomputation.
\item A calibrated error estimate on every correspondence.
\end{enumerate}

If an application consumes one full dense flow field per frame and nothing else,
v2 is the better choice on every axis: more accurate, and 11\% faster in that mode.
I would say so plainly. v3 earns its place when the consumer picks its own query
points, revisits a frame, needs positions between pixels, or wants a confidence
value, and when 2.6\% of mean accuracy is an acceptable price for those.

\section{Limitations}
\label{sec:limitations}

\textbf{Accuracy.} The decoder is less accurate than the upsampler it replaces, by
2.6\% on mean error and 1.5 points of 1-pixel accuracy at best, 7.6\% on the
like-for-like comparison. I know the cause: the decoder's finest input is at $1/8$
resolution, so the evidence it sees barely changes inside an $8\times8$ cell, while
v2's upsampler reads the full-resolution frame directly. Three independent
experiments converge on that same mechanism. A Fourier positional encoding of the
sub-cell offset returned null, which rules out a missing positional signal. Sub-pixel
querying gains very little. Querying below the input sampling rate matches bilinear
interpolation exactly. That convergence is what makes the resolution ceiling a
diagnosis rather than a suspicion. It is diagnosed and unfixed in this work.

The same limitation means sub-pixel querying is currently closer to interpolation
than to genuine sub-pixel inference: queries within one $8\times8$ cell see nearly
identical evidence. The interface is exact and the mechanism is real, but the extra
information it can return between pixels is small until the decoder sees
full-resolution features.

\textbf{No embedded measurement.} Every latency figure is from a laptop RTX 4060. The
edge claim is a design target, not a result. The region speedup is hardware
dependent for the same reason: a small region becomes launch-overhead bound on a
larger accelerator. The accuracy cost of a region is not hardware dependent.

\textbf{One primary evaluation domain.} VKITTI2 is synthetic driving footage, and
it says little on its own about field or survey imagery. The 40 aerial TartanAir
pairs answer that only partly.

\textbf{Statistical resolution.} One seed per configuration, with
checkpoint-to-checkpoint variation up to 0.038\px. Differences below about
0.05\px\ are not resolved.

\section{Next}
\label{sec:next}
Ordered by what each item settles, not by how much work it takes.

\begin{enumerate}
\item \textbf{Embedded measurement on a Jetson.} Converts the edge claim into a
      result, or refutes it. Nothing else in this work depends on the measurement
      platform the way this does.
\item \textbf{A second seed per configuration}, one training run each. Three of
      the four configurations currently sit inside the noise floor and cannot be
      ranked. A second seed is what would allow ranking them honestly.
\item \textbf{Evaluation above the input sampling rate.} Query at $2\times$
      resolution against native 4K ground truth. The script exists and runs
      today. This is the one test v2 structurally cannot take, because the
      question cannot be posed to it.
\item \textbf{Tiled refinement.} Refinement is 60\% of runtime and all of its
      operations are local, so it should run on region tiles with a halo, and the
      measured halo decay indicates no retraining is needed. Estimated
      2$\times$. This is not a v3 contribution: it would benefit v2 equally.
\item \textbf{Registration on field and survey imagery}, which tests the intended
      application and addresses the single-domain limitation directly.
\end{enumerate}

\begingroup
\small
\begin{thebibliography}{99}

\bibitem{hs81}
B.~K.~P. Horn and B.~G. Schunck, ``Determining optical flow,'' \emph{Artificial
Intelligence}, vol.~17, no.~1--3, pp.~185--203, 1981.

\bibitem{lk81}
B.~D. Lucas and T.~Kanade, ``An iterative image registration technique with an
application to stereo vision,'' in \emph{Proc. Int. Joint Conf. Artificial
Intelligence (IJCAI)}, 1981, pp.~674--679.

\bibitem{brox04}
T.~Brox, A.~Bruhn, N.~Papenberg, and J.~Weickert, ``High accuracy optical flow
estimation based on a theory for warping,'' in \emph{Proc. European Conf.
Computer Vision (ECCV)}, 2004.

\bibitem{flownet}
A.~Dosovitskiy, P.~Fischer, E.~Ilg, P.~H\"{a}usser, C.~Hazirbas, V.~Golkov,
P.~van~der~Smagt, D.~Cremers, and T.~Brox, ``FlowNet: Learning optical flow with
convolutional networks,'' in \emph{Proc. IEEE Int. Conf. Computer Vision (ICCV)},
2015.

\bibitem{flownet2}
E.~Ilg, N.~Mayer, T.~Saikia, M.~Keuper, A.~Dosovitskiy, and T.~Brox, ``FlowNet
2.0: Evolution of optical flow estimation with deep networks,'' in \emph{Proc.
IEEE Conf. Computer Vision and Pattern Recognition (CVPR)}, 2017.

\bibitem{pwcnet}
D.~Sun, X.~Yang, M.-Y. Liu, and J.~Kautz, ``PWC-Net: CNNs for optical flow using
pyramid, warping, and cost volume,'' in \emph{Proc. IEEE Conf. Computer Vision
and Pattern Recognition (CVPR)}, 2018.

\bibitem{liteflownet}
T.-W. Hui, X.~Tang, and C.~C. Loy, ``LiteFlowNet: A lightweight convolutional
neural network for optical flow estimation,'' in \emph{Proc. IEEE Conf. Computer
Vision and Pattern Recognition (CVPR)}, 2018.

\bibitem{raft}
Z.~Teed and J.~Deng, ``RAFT: Recurrent all-pairs field transforms for optical
flow,'' in \emph{Proc. European Conf. Computer Vision (ECCV)}, 2020.

\bibitem{gmflow}
H.~Xu, J.~Zhang, J.~Cai, H.~Rezatofighi, and D.~Tao, ``GMFlow: Learning optical
flow via global matching,'' in \emph{Proc. IEEE/CVF Conf. Computer Vision and
Pattern Recognition (CVPR)}, 2022.

\bibitem{neuflow1}
Z.~Zhang, H.~Jiang, and H.~Singh, ``NeuFlow: Real-time, high-accuracy optical
flow estimation on robots using edge devices,'' \emph{arXiv:2403.10425}, 2024.

\bibitem{neuflow2}
Z.~Zhang, A.~Gupta, H.~Jiang, and H.~Singh, ``NeuFlow v2: High-efficiency optical
flow estimation on edge devices,'' \emph{arXiv:2408.10161}, 2024.

\bibitem{nerf}
B.~Mildenhall, P.~P. Srinivasan, M.~Tancik, J.~T. Barron, R.~Ramamoorthi, and
R.~Ng, ``NeRF: Representing scenes as neural radiance fields for view
synthesis,'' in \emph{Proc. European Conf. Computer Vision (ECCV)}, 2020.

\bibitem{siren}
V.~Sitzmann, J.~N.~P. Martel, A.~W. Bergman, D.~B. Lindell, and G.~Wetzstein,
``Implicit neural representations with periodic activation functions,'' in
\emph{Advances in Neural Information Processing Systems (NeurIPS)}, 2020.

\bibitem{liif}
Y.~Chen, S.~Liu, and X.~Wang, ``Learning continuous image representation with
local implicit image function,'' in \emph{Proc. IEEE/CVF Conf. Computer Vision
and Pattern Recognition (CVPR)}, 2021.

\bibitem{anyflow}
H.~Jung, Z.~Hui, L.~Luo, H.~Yang, F.~Liu, S.~Yoo, R.~Ranjan, and D.~Demandolx,
``AnyFlow: Arbitrary scale optical flow with implicit neural representation,'' in
\emph{Proc. IEEE/CVF Conf. Computer Vision and Pattern Recognition (CVPR)}, 2023.

\bibitem{infinidepth}
H.~Yu, H.~Lin, J.~Wang, J.~Li, Y.~Wang, X.~Zhang, Y.~Wang, X.~Zhou, R.~Hu, and
S.~Peng, ``InfiniDepth: Arbitrary-resolution and fine-grained depth estimation
with neural implicit fields,'' in \emph{Proc. IEEE/CVF Conf. Computer Vision and
Pattern Recognition (CVPR)}, 2026. arXiv:2601.03252.

\bibitem{kendall}
A.~Kendall and Y.~Gal, ``What uncertainties do we need in Bayesian deep learning
for computer vision?'' in \emph{Advances in Neural Information Processing Systems
(NeurIPS)}, 2017.

\bibitem{ilg18}
E.~Ilg, O.~Cicek, S.~Galesso, A.~Klein, O.~Makansi, F.~Hutter, and T.~Brox,
``Uncertainty estimates and multi-hypotheses networks for optical flow,'' in
\emph{Proc. European Conf. Computer Vision (ECCV)}, 2018.

\bibitem{pdcnet}
P.~Truong, M.~Danelljan, L.~Van~Gool, and R.~Timofte, ``Learning accurate dense
correspondences and when to trust them,'' in \emph{Proc. IEEE/CVF Conf. Computer
Vision and Pattern Recognition (CVPR)}, 2021.

\end{thebibliography}
\endgroup

\vspace{6pt}
{\color{rule}\hrule}
\vspace{4pt}
{\footnotesize\color{muted}\raggedright
Code, training configurations and the full development record:
\texttt{github.com/shrirag10/Neuflowv3}. Every number here is reproducible from
\texttt{scripts/eval\_all\_runs.py}, \texttt{scripts/benchmark\_sparse.py},
\texttt{scripts/eval\_calibration.py}, \texttt{scripts/eval\_roi\_crop.py} and
\texttt{scripts/bench\_scenarios.py}. Full parameter provenance is in
\texttt{docs/base\_parameters.md}.\par}

\end{document}
